% GENERATED FILE. Edit canonical lessons, then run npm run generate:books.
% canonical-source-hash: fnv1a64:e67c1a3c5d19957a
% canonical-lessons: RU-C311-prazdnovat, RU-C311-otmechat, RU-C311-torzhestvo, RU-C311-torzhestvennyy, RU-C311-prazdnichnyy

\chapter{Celebrating}
\label{ch:ru-celebrating}
\hlchaptermodality{311}

\begin{chapteropening}
\textbf{By the end of this chapter:} \emph{I can say that people celebrated or marked an occasion, and talk about a celebration, a solemn moment and a festive meal.}

Complete the last lesson of chapter 311: I can say that people celebrated or marked an occasion, and talk about a celebration, a solemn moment and a festive meal.
\end{chapteropening}

\section[\texorpdfstring{\ru{праздновать} (prázdnovat)}{праздновать (prázdnovat)}]{\ru{праздновать} (prázdnovat) — to celebrate}

\label{lesson:RU-C311-prazdnovat}



\begin{quote}
Before the new one: say the Russian for a christening, then the Russian for to christen.
\end{quote}

\subsection*{You'll want to know: \ru{праздновать}}
\textbf{\ru{праздновать}} (\emph{prázdnovat}) — \textquotedblleft{}to celebrate\textquotedblright{}.

\begin{quote}
\ru{Мы} \ru{праздновали} \ru{день} \ru{рождения} \ru{до} \ru{поздней} \ru{ночи}.
\end{quote}

\begin{quote}
\emph{My prázdnovali den rozhdéniya do pózdney nóchi.}
\end{quote}

\begin{quote}
We celebrated the birthday until late at night.
\end{quote}

Imperfective; the perfective is \emph{\ru{отпраздновать}} (\emph{otprázdnovat}). Present \emph{\ru{праздную}}.

\begin{cousinweb}
\emph{\ru{праздновать}} is built on \emph{\ru{праздник}} \textquotedblleft{}a holiday\textquotedblright{} (EtymDB), from \emph{\ru{праздный}} \textquotedblleft{}idle\textquotedblright{}, the Church Slavonic shape of native \emph{\ru{порожний}} \textquotedblleft{}empty\textquotedblright{}: a day free of work.
\end{cousinweb}

\begin{grammarlens}[title={What you've built}]
One more word for telling, explaining and describing at B1.
\end{grammarlens}

\subsection*{Your turn}
\begin{itemize}
\raggedright
  \item \emph{Say it:} \emph{prázdnovat}
  \item \emph{Say it:} \emph{prázdnovat}, once more
  \item \emph{Say it:} a sentence of your own with \emph{prázdnovat}
  \item \emph{Recall:} say the Russian for a christening, then the Russian for to christen, then say \emph{prázdnovat} again
\end{itemize}

\begin{tcolorbox}[breakable,colback=teal!4,colframe=teal!35!black,title={Before you move on}]
What does \ru{праздновать} mean? (\textquotedblleft{}To celebrate\textquotedblright{}.) Say it once more.
\end{tcolorbox}
\section[\texorpdfstring{\ru{отмечать} (otmechát)}{отмечать (otmechát)}]{\ru{отмечать} (otmechát) — to mark (an occasion), to celebrate; to note}

\label{lesson:RU-C311-otmechat}



\begin{quote}
Before the new one: say the Russian for to christen, then the Russian for to celebrate.
\end{quote}

\subsection*{You'll want to know: \ru{отмечать}}
\textbf{\ru{отмечать}} (\emph{otmechát}) — \textquotedblleft{}to mark (an occasion), to celebrate; to note\textquotedblright{}.

\begin{quote}
\ru{Каждый} \ru{год} \ru{мы} \ru{отмечаем} \ru{этот} \ru{день} \ru{вместе}.
\end{quote}

\begin{quote}
\emph{Kázhdyy god my otmecháyem étot den vméste.}
\end{quote}

\begin{quote}
Every year we mark this day together.
\end{quote}

Imperfective; perfective \emph{\ru{отметить}} (\emph{otmétit}), \emph{\ru{отмечу}}. The everyday word for celebrating: \emph{\ru{отметить} \ru{день} \ru{рождения}}.

\begin{cousinweb}
\emph{\ru{отмечать}} is built on \emph{\ru{от}-} and \emph{\ru{метить}} \textquotedblleft{}to mark\textquotedblright{}, inherited from Proto-Slavic \emph{mětiti} (EtymDB): to set a mark on a day.
\end{cousinweb}

\begin{grammarlens}[title={What you've built}]
One more word for telling, explaining and describing at B1.
\end{grammarlens}

\subsection*{Your turn}
\begin{itemize}
\raggedright
  \item \emph{Say it:} \emph{otmechát}
  \item \emph{Say it:} \emph{otmechát}, once more
  \item \emph{Say it:} a sentence of your own with \emph{otmechát}
  \item \emph{Recall:} say the Russian for to christen, then the Russian for to celebrate, then say \emph{otmechát} again
\end{itemize}

\begin{tcolorbox}[breakable,colback=teal!4,colframe=teal!35!black,title={Before you move on}]
What does \ru{отмечать} mean? (\textquotedblleft{}To mark (an occasion), to celebrate; to note\textquotedblright{}.) Say it once more.
\end{tcolorbox}
\section[\texorpdfstring{\ru{торжество} (torzhestvó)}{торжество (torzhestvó)}]{\ru{торжество} (torzhestvó) — a celebration, a festive occasion; triumph}

\label{lesson:RU-C311-torzhestvo}



\begin{quote}
Before the new one: say the Russian for to celebrate, then the Russian for to mark.
\end{quote}

\subsection*{You'll want to know: \ru{торжество}}
\textbf{\ru{торжество}} (\emph{torzhestvó}) — \textquotedblleft{}a celebration, a festive occasion; triumph\textquotedblright{}.

\begin{quote}
\ru{Торжество} \ru{началось} \ru{в} \ru{шесть} \ru{часов}.
\end{quote}

\begin{quote}
\emph{Torzhestvó nachalós v shest chasóv.}
\end{quote}

\begin{quote}
The celebration began at six.
\end{quote}

Neuter.

\begin{cousinweb}
\emph{\ru{торжество}} is inherited from Old East Slavic \emph{tŭržĭstvo} (EtymDB), usually linked with \emph{\ru{торг}} \textquotedblleft{}a market, a gathering of people\textquotedblright{}.
\end{cousinweb}

\begin{grammarlens}[title={What you've built}]
One more word for telling, explaining and describing at B1.
\end{grammarlens}

\subsection*{Your turn}
\begin{itemize}
\raggedright
  \item \emph{Say it:} \emph{torzhestvó}
  \item \emph{Say it:} \emph{torzhestvó}, once more
  \item \emph{Say it:} a sentence of your own with \emph{torzhestvó}
  \item \emph{Recall:} say the Russian for to celebrate, then the Russian for to mark, then say \emph{torzhestvó} again
\end{itemize}

\begin{tcolorbox}[breakable,colback=teal!4,colframe=teal!35!black,title={Before you move on}]
What does \ru{торжество} mean? (\textquotedblleft{}A celebration, a festive occasion; triumph\textquotedblright{}.) Say it once more.
\end{tcolorbox}
\section[\texorpdfstring{\ru{торжественный} (torzhéstvennyy)}{торжественный (torzhéstvennyy)}]{\ru{торжественный} (torzhéstvennyy) — solemn, ceremonial}

\label{lesson:RU-C311-torzhestvennyy}



\begin{quote}
Before the new one: say the Russian for to mark, then the Russian for a celebration.
\end{quote}

\subsection*{You'll want to know: \ru{торжественный}}
\textbf{\ru{торжественный}} (\emph{torzhéstvennyy}) — \textquotedblleft{}solemn, ceremonial\textquotedblright{}.

\begin{quote}
\ru{Это} \ru{был} \ru{торжественный} \ru{момент}.
\end{quote}

\begin{quote}
\emph{Éto byl torzhéstvennyy momént.}
\end{quote}

\begin{quote}
It was a solemn moment.
\end{quote}

\emph{\ru{Торжественно}} — solemnly.

\begin{cousinweb}
\emph{\ru{торжественный}} is built on \emph{\ru{торжество}} \textquotedblleft{}a celebration\textquotedblright{}, inherited from Old East Slavic \emph{tŭržĭstvo} (EtymDB).
\end{cousinweb}

\begin{grammarlens}[title={What you've built}]
One more word for telling, explaining and describing at B1.
\end{grammarlens}

\subsection*{Your turn}
\begin{itemize}
\raggedright
  \item \emph{Say it:} \emph{torzhéstvennyy}
  \item \emph{Say it:} \emph{torzhéstvennyy}, once more
  \item \emph{Say it:} a sentence of your own with \emph{torzhéstvennyy}
  \item \emph{Recall:} say the Russian for to mark, then the Russian for a celebration, then say \emph{torzhéstvennyy} again
\end{itemize}

\begin{tcolorbox}[breakable,colback=teal!4,colframe=teal!35!black,title={Before you move on}]
What does \ru{торжественный} mean? (\textquotedblleft{}Solemn, ceremonial\textquotedblright{}.) Say it once more.
\end{tcolorbox}
\section[\texorpdfstring{\ru{праздничный} (prázdnichnyy)}{праздничный (prázdnichnyy)}]{\ru{праздничный} (prázdnichnyy) — festive, holiday}

\label{lesson:RU-C311-prazdnichnyy}



\begin{quote}
Before the new one: say the Russian for a celebration, then the Russian for solemn.
\end{quote}

\subsection*{You'll want to know: \ru{праздничный}}
\textbf{\ru{праздничный}} (\emph{prázdnichnyy}) — \textquotedblleft{}festive, holiday\textquotedblright{}.

\begin{quote}
\ru{Мать} \ru{приготовила} \ru{праздничный} \ru{ужин}.
\end{quote}

\begin{quote}
\emph{Mat prigotóvila prázdnichnyy úzhin.}
\end{quote}

\begin{quote}
Mother made a festive dinner.
\end{quote}

\emph{\ru{Праздничный} \ru{стол}} — a holiday table.

\begin{cousinweb}
\emph{\ru{праздничный}} is built on \emph{\ru{праздник}} \textquotedblleft{}a holiday\textquotedblright{} (EtymDB), from \emph{\ru{праздный}} \textquotedblleft{}idle\textquotedblright{}, the Church Slavonic shape of native \emph{\ru{порожний}} \textquotedblleft{}empty\textquotedblright{}.
\end{cousinweb}

\begin{grammarlens}[title={What you've built}]
One more word for telling, explaining and describing at B1.
\end{grammarlens}

\subsection*{Your turn}
\begin{itemize}
\raggedright
  \item \emph{Say it:} \emph{prázdnichnyy}
  \item \emph{Say it:} \emph{prázdnichnyy}, once more
  \item \emph{Say it:} a sentence of your own with \emph{prázdnichnyy}
  \item \emph{Recall:} say the Russian for a celebration, then the Russian for solemn, then say \emph{prázdnichnyy} again
\end{itemize}

\begin{tcolorbox}[breakable,colback=teal!4,colframe=teal!35!black,title={Before you move on}]
What does \ru{праздничный} mean? (\textquotedblleft{}Festive, holiday\textquotedblright{}.) Say it once more.
\end{tcolorbox}
